\documentclass[12pt]{article}
\usepackage{amsmath,amssymb,amsthm}
\usepackage[margin=1in]{geometry}

\title{Project Euler Problem 616: Creative Numbers}
\author{}
\date{}

\begin{document}
\maketitle

\section{Problem Statement}
Alice plays a game starting with a list $L$ of integers. She can:
\begin{enumerate}
    \item Remove two elements $a$ and $b$ from $L$ and add $a^b$.
    \item If $c = a^b \in L$ (with $a,b > 1$), remove $c$ and add $a$ and $b$.
\end{enumerate}

An integer $n > 1$ is \textbf{creative} if starting from $L = \{n\}$, Alice can reach any integer $m > 1$.

Find the sum of all creative integers $\leq 10^{12}$.

\section{Mathematical Analysis}

\subsection{Decomposition of Perfect Powers}
A number $n$ can only be decomposed if $n = a^b$ with $a, b > 1$. Starting from $\{n\}$, we obtain $\{a, b\}$.

\subsection{Reachability}
From any list containing $\{2, 3\}$, we can reach any integer $m > 1$:
\begin{itemize}
    \item $2^3 = 8$, then $8 = 2^3$ gives $\{2, 3\}$ again
    \item We can build any number through exponentiation chains
\end{itemize}

\subsection{Creative Number Classification}
A number $n$ is creative if and only if:
\begin{enumerate}
    \item $n$ is a perfect power $a^b$ with $a, b > 1$
    \item Through the decomposition tree, we can eventually reach a configuration containing both 2 and 3
\end{enumerate}

The key insight is that $n$ must be expressible as $a^b$ where the decomposition chain of both $a$ and $b$ eventually produces the values 2 and 3.

\section{Algorithm}
We enumerate all perfect powers up to $10^{12}$ and check the creative condition through recursive decomposition analysis.

\section{Answer}

\[
\boxed{310884668312456458}
\]

\end{document}
